\ifthenelse{\equal{\lang}{ko}}{%
우리의 결과는 자동화된 규제 감지의 가능성을 보여주지만, 여전히 몇 가지 한계가 남아 있다. 
~\\
\textbf{언어 커버리지:} 현재 평가는 영어, 한국어, 프랑스어에 한정되어 있으며, 
저자원 언어 확장을 위해서는 추가 데이터와 다국어 적응이 필요하다. 
스페인어, 아랍어, 베트남어에 대한 예비 실험에서는 제로샷 성능이 0.71 macro-F1이었으나, 
200개의 소수샷 예제를 활용할 경우 0.84로 향상되었다. 
이는 본 분류체계와 파이프라인이 최소한의 감독으로도 글로벌 적응이 가능함을 시사한다.  
~\\
\textbf{개념 드리프트:} 법률 및 정책 환경의 변화는 시간이 지남에 따라 정확도를 감소시킬 수 있다. 
3개월 롤링 평가에서 F1=0.89에서 0.74로 하락했으나, 
월별 경량 재학습(3 에폭, 5천 샘플)을 통해 0.88로 회복되었다. 
향후 연구에서는 자동 드리프트 탐지 및 연속 학습을 통해 지속적인 강건성을 확보할 것이다.  
\textbf{인간 참여 확장성:} 분석가 피드백은 정확도를 향상시키지만 법률 전문가 의존성이 크므로, 
소규모 조직을 위해서는 경량화된 프로토콜이 필요하다.  
\textbf{배포 비용:} QLoRA는 연산 자원 소모를 줄이지만, 
대규모 활용 시 GDPR 환경에서 비용 및 프라이버시 문제가 여전히 존재한다. 
~\\
\textbf{방법론적 혁신(Methodological Innovation):} 
본 연구는 기존 LLM 및 NLP 기법을 성공적으로 통합하여 규제 모니터링 프레임워크를 구축하였으나, 
향후 연구에서는 \textit{알고리즘적 독창성}을 강화할 필요가 있다. 
예를 들어, 규제 문서에 특화된 새로운 프롬프트 엔지니어링 기법을 개발하거나, 
거버넌스 과제에 최적화된 경량 모델 아키텍처를 설계하는 방법이 있다. 
또한 인간 전문가 피드백(RLHF 유사 메커니즘)의 정량적 효과를 측정하는 통제 실험을 포함한다면, 
제안된 파이프라인의 효과성과 신뢰성에 대한 더 강력한 경험적 근거를 제시할 수 있을 것이다. 
이러한 방법론적 혁신은 LexSense를 고도로 실용적인 시스템에서, 
학술적으로도 독창성이 높은 기여로 발전시킬 것이다. 

종합적으로, 이러한 한계와 확장은 전 세계적 규모의 AI 거버넌스 정보학으로 나아가는 경로를 제시하며, 
효율성, 강건성, 그리고 공정성·투명성·프라이버시 원칙과의 정렬을 보장한다.


}{%

While our results demonstrate the promise of automated regulatory sensing, several limitations remain. 
~\\
\textbf{Language coverage:} current evaluation is limited to English, Korean, and French; extending to low-resource languages requires additional data and multilingual adaptation. Preliminary experiments on Spanish, Arabic, and Vietnamese showed zero-shot performance of 0.71 macro-F1, which improved to 0.84 with 200 few-shot examples, indicating that our taxonomy and pipeline can adapt globally with minimal supervision.
~\\
\textbf{Concept drift:} evolving legal and policy landscapes may reduce accuracy over time. In a three-month rolling evaluation, performance dropped from F1=0.89 to 0.74, but lightweight monthly retraining (3 epochs, 5k samples) restored it to 0.88. Future work will explore automated drift detection and continual learning for sustainable robustness. Human-in-the-loop scalability: analyst feedback improves accuracy but depends on legal experts, so lightweight protocols are needed for smaller organizations. Deployment costs: although QLoRA reduces compute demands, large-scale use still faces cost and privacy challenges under GDPR. 
~\\
\textbf{Methodological Innovation:} While this work successfully integrates existing LLM and NLP techniques into a unified regulatory monitoring framework, future research should aim to strengthen its \textit{algorithmic originality}. For example, developing novel prompt engineering strategies tailored to regulatory documents or designing lightweight model architectures specialized for governance tasks could go beyond the current integration-based approach. Moreover, including controlled experiments that quantitatively measure the impact of human expert feedback (RLHF-style mechanisms) would provide stronger empirical evidence of the effectiveness and reliability of the proposed pipeline. Such methodological innovations would further elevate LexSense from a highly practical system to a contribution of the highest academic originality. Together, these limitations and extensions illustrate a path toward AI Governance Informatics at global scale, ensuring efficiency, robustness, and alignment with fairness, transparency, and privacy principles. 
}